\chapter{INTRODUCTION}

\section{SIGNIFICANCE AND BACKGROUND OF THE PROBLEM}
\subsection{Neglecting Seabed Pollution}

Thailand’s marine territories, for example, the Andaman Sea and the Gulf of
Thailand, are becoming more polluted. A significant amount of pollution, including plastics, heavy metals, industrial runoff, and abandoned fishing equipment, falls to the seabed. Which is affecting organisms and disrupting the food chain. In contrast to surface waste, which is frequently noticeable and easily removed, seabed contamination is concealed and keeps building up. The problem is compounded by the lack of inspection, creating blind spots in ecosystem management and marine policy.

This calls for an innovative solution: a system that can continuously and
independently monitor undersea conditions. UUVs are transformative in this situation. They provide the capacity to explore the ocean floor, identify environmental risks, and gather information that can guide prompt action.

\subsection{Enhancing Marine Species Monitoring in Thailand}

Monitoring pollution levels and marine biodiversity is crucial for managing
ecosystems sustainably. As markers of the health of the environment, marine species' loss frequently indicates more serious problems with habitat degradation, temperature changes, or water quality.

However, traditional maritime monitoring techniques that are commonly
used in Thailand, such as diver-based surveys or equipment installed on boats, are risky, time-consuming, and labor-intensive. These techniques are further constrained by depth limitations, strong currents, and low vision.

\subsection{Ensuring Maintenance of Underwater Systems}

Underwater systems, for example, pipelines, fish farms, and marine research
equipment, are constantly exposed to harsh environmental conditions including corrosion, and damage from marine life. These systems require regular inspection, repair, and maintenance to remain functional and safe.


Traditional inspection methods rely on divers or expensive Remotely
Operated Vehicles (ROVs), which are costly, time-consuming, and pose safety risks especially in deep or murky waters. As a result, many installations suffer from delayed maintenance and undetected failures, leading to higher repair costs or system downtime.

\subsection{Traditional UUV is costly}

Unmanned Underwater Vehicles are essential for tasks such as marine
research, environmental monitoring, and infrastructure inspection. However, most traditional UUVs are expensive and inaccessible to small research labs, universities, and local conservation groups.For instance, Chasing gladius Mini (35,000 –60,000 Thai baht), BlueROV2 (180,000 Thai baht), and VideoRay Pro 5 (500,000 – 3,700,000 Thai baht). These systems also require additional costs for operation, deployment, and training. Due to these barriers, many underwater systems go uninspected, and marine research is limited—especially in developing regions like Thailand.


\section{OBJECTIVES OF THE PROJECT}

The primary objective of this project is to design and develop multipurpose
Unmanned Underwater Vehicle (UUV) that supports various underwater applications in both research and maintenance. The system is aimed to provide affordable, modular, and adaptable platform that can operate effectively in coastal and nearshore environments of Thailand.\\
\textbf{Key Objectives}
\begin{enumerate}
	\item Develop a Modular Multi-Use UUV Platform:	Create a compact and cost-effective UUV structure that allows easy integration of various tools and sensors based on mission requirements. Design the system to support plug-and-play modules for quick reconfiguration.
	\item Enable Underwater Debris Detection and Cleaning:
	Integrate vision-based detection for identifying marine litter and seabed debris. Attach basic cleaning tools or manipulators for collecting or tagging debris for environmental maintenance.
	\item Support Marine Biological Research (Ichthyology):
	Equip the UUV with high-resolution cameras to monitor marine life habitats and behaviors. Enable data collection for research in ichthyology and marine biodiversity mapping.
	\item Monitor and Inspect Underwater Infrastructure:
	Facilitate routine inspection of underwater systems such as pipelines, fish cages, research sensors, and marine installations. Provide real-time visual data for structural assessment and potential failure detection.
\end{enumerate}

\section{SCOPE AND LIMITATIONS OF PROJECT}

To minimize manufacturing costs while maintaining structural functionality, the connecting components of the Unmanned Underwater Vehicle (UUV) were fabricated by 3d-printing PETG components before applying a resin shell around. This method allows for rapid prototyping, ease of customization, and acceptable water resistance, particularly in shallow-to-moderate depth applications. Traditional 3D-printed materials such as PLA, while accessible, exhibit poor waterproofing and dimensional stability in underwater environments, making them unsuitable for critical structural components.